% !TeX root = ../../main.tex
\subsection{Normal convergence away from a critical point}
\label{subsec:b:regular}

At a simple inverse branch, the coefficients $L_{g,h}$ have an
unambiguous meaning before their Borel series is known to converge:
evaluate the rational jet expressions of
Lemma~\ref{lmm:s07:finite-jets} on that branch. Their characteristic
identities remain identities of those rational expressions. We now
use them to construct the actual holomorphic Borel functions.

The first estimate is local and has no critical-point hypothesis. Its
dependence on the inverse derivative is important: this derivative
grows along a remote characteristic and on the negative real axis.

\begin{lmm}[A regular disk]\label{lmm:b:regular-disk}
Let $u_0\in\mathbb C$, $\theta_0\geq0$, $r>0$ and $\beta>0$,
and suppose that
\begin{equation}\label{eq:b:regular-hypothesis}
 |w_{\theta_0}'(u)-\beta|\leq\beta/4
                  \quad\text{on }\overline{\mathbb D(u_0,r)}.
\end{equation}
For every fixed integer $h\geq1$ there are positive constants
$c_h,C_h$, depending only on $h$, such that, for every
$\theta\geq\theta_0$, the Borel series $\mathcal L_h$ on this
inverse branch converges normally when
\begin{equation}\label{eq:b:regular-disk}
 |x-w_\theta(u_0)|+|y|
                  <c_h r(\beta+\theta-\theta_0).
\end{equation}
Its supremum there is at most
$C_h(\beta+\theta-\theta_0)^{h-1}$. On a smaller disk, for every
fixed pair of nonnegative integers $a,k$,
\begin{equation}\label{eq:b:regular-derivatives}
 |\partial_x^a\partial_y^k\mathcal L_h|
 \leq C_{h,a,k}(\beta+\theta-\theta_0)^{h-1}
                  [r(\beta+\theta-\theta_0)]^{-a-k}.
\end{equation}
The characteristic integrals defining these functions converge
normally on smaller disks. The constants are independent of
$u_0,r,\beta,\theta_0$ and of the genus.
\end{lmm}
\begin{prf}
Write $s=\beta+\theta-\theta_0>0$ only within this proof.
Condition~\eqref{eq:b:regular-hypothesis} bounds the derivative of
$w_\theta$ above and below by fixed multiples of $s$, with
positive real part. Radial integration and univalence give an inverse
disk of radius a fixed multiple of $r s$. On a smaller
disk, its derivative differs from its central value by a small fixed
relative amount, by Cauchy's estimate. The normalized difference
quotient in \eqref{eq:s07:logarithmic-kernel} is therefore bounded
away from zero and has a bounded logarithm. Its odd difference at
$y/m$ is $O(|y|/(m r s))$. Equation~\eqref{eq:s07:borel-one}
and $\sum m^{-2}<\infty$ prove the assertion for $h=1$.

Suppose the assertion holds below $h$. Choose the radius constant
for layer $h$ smaller than fixed multiples of those lower-layer
constants. Along a characteristic starting in this disk,
\[
 |v_\theta(x)-u_0|
       \leq C|x-w_\theta(u_0)|/s,
\]
so its displacement from $w_{\theta+\sigma}(u_0)$ is at most
$C|x-w_\theta(u_0)|(s+\sigma)/s$.
All shifted lower-layer arguments thus remain in their smaller disks.

The cubic shift remainder in $\mathcal K$ is bounded by the third
$x$ derivative times $|(1-\tau)y/m|^3$; its outer factor $m$
leaves a convergent $m^{-2}$ sum. The same estimate after one
$y$ derivative bounds $1+y\partial_y$. Its normalized disk norm
is $O((s+\sigma)^{h-2})$. In $\mathcal N_h$, the two
$x$ derivatives cost $(r (s+\sigma))^{-2}$ and $y^2$
restores those two powers. The lower-layer norms multiply to
$O((s+\sigma)^{h-2})$ as well.

Both sources vanish to order at least $3h-2$ in $y$, by the
coefficient calculation in Proposition~\ref{prp:s07:borel-transport}.
Cauchy's estimate on two nested disks therefore gives, on the
starting smaller disk, the integrable majorant
\begin{equation}\label{eq:b:regular-integral}
 C_h (s+\sigma)^{h-2}
         \left(\frac{s}{(s+\sigma)}\right)^{3h-2}.
\end{equation}
Its integral over $\sigma\geq0$ is
$C_h s^{h-1}/(2h-1)$. The characteristic formula now
defines a holomorphic function with the asserted norm. Normal
convergence allows coefficient extraction; the coefficients are
exactly those of \eqref{eq:s07:borel-transport}, so this function
is the required Borel sum. This proves the induction.
Ordinary Cauchy estimates on smaller disks give
\eqref{eq:b:regular-derivatives}, including any fixed derivative
orders. No infinite-genus interchange is assumed before this
integrable majorant has been established.
\end{prf}

For a bounded convex domain $\Omega$ with
$\operatorname{Re}w_{\theta_0}'>0$ in $\Omega$, and
$\theta\geq\theta_0$, define the open tube
\begin{equation}\label{eq:b:tube}
 \mathfrak T_\theta^\Omega
 =\{(x,y):x\in w_\theta(\Omega),\quad
       |y|<\operatorname{dist}(x,\partial w_\theta(\Omega))\}.
\end{equation}
Assume $w$ is holomorphic on a neighborhood of the closure; the
domains used here are disks, so this causes no boundary regularity
issue. Compact subsets of these tubes are the convenient bulk
domains for the kernel integrals.

\begin{prp}[Holomorphic transport on inverse tubes]\label{prp:b:tube}
Every $\mathcal L_h$, $h\geq1$, converges normally and is
holomorphic on $\mathfrak T_\theta^\Omega$. Its characteristic
integral and the scalar sums in the kernels converge normally on
compact subsets. For the domains $D_\theta$ of
Proposition~\ref{prp:b:inverse-budget}, and real $0\leq x<R_\theta$
(also on the small negative interval admitted there),
the function also extends jointly in $(x,y)$ across every point of
$|y|=R_\theta-x$ other than $y=\pm(R_\theta-x)$.
These continuation neighborhoods can be chosen uniformly when
$\theta,x$ and the noncritical normalized arc range over compact
subsets of their stated ranges.
\end{prp}
\begin{prf}
The inverse is univalent, and its derivative has positive real part.
For $h=1$ its normalized difference quotient has no zero for
$|y|$ less than the inverse radius. It has a holomorphic logarithm
obtained by continuation from $y=0$; equivalently, its nonvanishing
argument is defined on a disk. The odd difference gives a normally
convergent $m^{-2}$ majorant on compact subsets.

Induct on $h$. A kernel shift satisfies $|c|+t\leq1$ and thus
sends a point of the tube strictly into the corresponding
lower-layer tube, by the triangle inequality for boundary distance.
Characteristic transport does not decrease that distance, by the
last argument of Lemma~\ref{lmm:b:cone-transport}, which applies
unchanged to this convex $\Omega$. On a bounded interval of
$\sigma$, the shifted arguments range over compact subsets of the
lower-layer tubes. Their holomorphic norms and the cubic-difference
$m^{-2}$ majorant are consequently uniform.

For large $\sigma$, the inverse points remain in a compact subset
of $\Omega$. Choose small disks about those points. After increasing
the base parameter, the derivative on each such disk differs from
that positive base parameter by at most one quarter of it. Apply
Lemma~\ref{lmm:b:regular-disk}. Its radius grows linearly with
$\sigma$, while the Borel argument stays bounded. The vanishing
order $3h-2$ gives the same integrable bound as
\eqref{eq:b:regular-integral}. Finitely many disks cover the compact
inverse set. Thus the whole characteristic integral converges
normally, and coefficient extraction identifies it with the formal
series. This proves the first assertion without an assumed analytic
Borel sum.

For continuation across the indicated circle, first observe that
the segment $\{x+c y:-1\leq c\leq1\}$ lies in $D_\theta$
when $|y|=R_\theta-x$ and $y\ne\pm(R_\theta-x)$.
This follows from the unique boundary contact in
\eqref{eq:b:inverse-budget}. For $h=1$, both inverse arguments
are interior points of $D_\theta$, and injectivity keeps their
difference quotient nonzero; its logarithm therefore continues
locally. For a higher layer, at $\sigma=0$ every shifted
lower-layer disk has a strict radius budget, except for the
unshifted endpoint at $t=1$. That endpoint is already continued
by induction on $h$. At $t=0$ its center is an interior point of
$D_\theta$ and its Borel variable is zero. The same reasoning
applies to the two quadratic factors; their possible $t\tau=1$
or $t(1-\tau)=1$ endpoint is again a lower layer.

These facts provide open neighborhoods for every point of the compact
finite-parameter integration set at $\sigma=0$. A finite subcover
also works for small $\sigma$. For $\sigma$ bounded away from zero,
the boundary-distance increase is strict, uniformly on compact
sets; the original tube argument applies. Large $\sigma$ is
controlled by the preceding regular-disk bound. Cauchy estimates on
these neighborhoods justify the normally convergent $m$ sum and
the joint holomorphic continuation. Parameter compactness proves
the final uniformity assertion.
\end{prf}

\begin{crl}[Remote characteristics]\label{crl:b:remote}
Fix a sufficiently small $\sigma_0>0$ and a finite layer $h\geq2$.
For $\theta$ in a small compact interval and $v$ near
$x_{*,\theta}$, the part of the layer-$h$ characteristic integral
over $\sigma\geq\sigma_0$ is holomorphic in $(v,y)$ for $y$ in
a fixed disk, and is $O_h(|y|^{3h-2})$. The same conclusion holds
after any fixed number of derivatives on smaller disks.
For real $x$ in a compact subinterval of $[0,R_\theta)$, that
remote integral is also holomorphic in a neighborhood of
$y=\pm(R_\theta-x)$, uniformly on such compact sets.
\end{crl}
\begin{prf}
At $v=x_{*,\theta}$ the derivative of $w_{\theta+\sigma_0}$ is
$\sigma_0$. On a fixed small disk about this point it differs
from $\sigma_0$ by at most $\sigma_0/4$, uniformly in small
$\theta$. The same holds for nearby inverse centers.
Lemma~\ref{lmm:b:regular-disk}, applied from this base parameter,
gives a radius comparable to $\sigma$ and the integrable source
bound $C_h|y|^{3h-2}\sigma^{-2h}$. Normal convergence and Cauchy
estimates prove the first assertion.
For centers in the stated compact real interval, transport through
$\sigma_0$ strictly increases the inverse-radius budget. The tube
proposition therefore gives a common neighborhood of both endpoints;
the same remote majorant controls its infinite end. No critical point
at a large value of $\theta+\sigma$ is needed in either argument.
\end{prf}
